\documentclass{article}
\usepackage{anyfontsize}
\usepackage{changepage}
\usepackage[T1]{fontenc} 
\usepackage[polish]{babel}
\usepackage[utf8]{inputenc}
\usepackage{array}
\usepackage{booktabs}
\usepackage{graphicx, float}
\usepackage{subcaption}
\usepackage[margin = 1.5cm]{geometry}
\usepackage{changepage}

\title{%
    Modelowanie układów przepływowych  \\
    \large Jak kształt batymetrii "W" wpływa na przestrzenny rozkład
    maksymalnych wysokości fali?}
\author{Adam Marszalik}
\date{}
\begin{document}

\maketitle

\section{Założenia oraz parametry symulacji}
 Celem tego projektu jest zbadanie jak batymetria w kształcie litery "W" wpływa na maksymalne wychylenie propagującej się fali.
 Równania płytkiej wody źle odwzoruwują ostre i strome krawędzie dlatego barymetrię wygładzono w sposób pokazany na poniższym rysunku.
 
 \begin{figure}[H]
  \centering
  \includegraphics[width=7cm]{./profil_batymetrii.pdf}
  \caption{Kształ dna zbiornika wodnego wzdłuż propagującej się fali.  }
  \label{fig:batymetria}
\end{figure}

Symulację wykonano na siatce o rozmiarze 40m na 170m, o kroku przestrzennym 1m. Krok czasowy wynosi 0.05sekundy, a całość symulacji jest ustawiona na 40 sekund. 
Grawitację ustawiono na $10 \frac{m}{s^2}$.Fala powstaje wzdłuż osi X na Y = 0 metrów z słupa wody wynoszącego 0.075m i szerokości 4 metrów.
Kluczowe w tym projekcie jest porównanie maksymalnej wysokości fali w miejscu pierwszego i drugiego wzniesiena na batymetrii.
Zbadanie różnicy w wysokości w tych miejscach pozwoli na określenie czy taki kształt powoduje zmiejszenie amplitudy fali.



\section{Wyniki symulacji}

 
  

\begin{figure}[H]
  \centering
  \begin{subfigure}{0.5\textwidth}
    \centering
    \includegraphics[width=1\linewidth]{./maksymalne_wychylenie_bezwzgledne.pdf}
    \caption{wykres}
    \label{fig:sub2_1}
    \end{subfigure}%
  \begin{subfigure}{0.5\textwidth}
    \centering
    \includegraphics[width=1\linewidth]{./mapa_maksymalnych_wychylen.pdf}
    \caption{mapa}
    \label{fig:sub2_2}
  \end{subfigure}
  \caption{Przestrzenny rozkład maksymalnych wysokości fali, wzdłuż  osi y dla x = 20 (a), mapa (b)}
  \label{fig:rozkład_1}
\end{figure}

  
  
  Na powyższych rysunkach widać charakterystyczne wzniesienia. Fala na drugim wzniesieniu widocznie się zmiejszyła. 
  Maksymalna wysokość na pierwszym wzniesieniu wynosi 0.050437 metrów, na drugim 0.040397 m. Różnica względla to 19.91\%.
  Na rysunku 2.(b) widać zmiejszenie się wychylenia fali wychodzące z krańców siatki. Najprawdopodobniej wynikają one z błednego warunku począdkowego słupa wody.
  W celu sprawdzenia poprawności przebiegu symulacji, wykonano mymulację o dwukrotnie zwiększonej rozdzielczości czasowej i przestrzennej. 
  Krok przestrzenny ustawiono na 0.5 metra a krok czasowy na 0.025s. Wyniki znajdują się na rysunkach poniżej.
  
  
  \begin{figure}[H]
    \centering
    \begin{subfigure}{0.5\textwidth}
      \centering
      \includegraphics[width=1\linewidth]{./maksymalne_wychylenie_podwojona_rozdzielczosc.pdf}
      \caption{Wykres}
      \label{fig:sub3_1}
      \end{subfigure}%
    \begin{subfigure}{0.5\textwidth}
      \centering
      \includegraphics[width=1\linewidth]{./mapa_maksymalnych_wychylen_podwojona.pdf}
      \caption{mapa}
      \label{fig:sub3_2}
    \end{subfigure}
    \caption{Przestrzenny rozkład maksymalnych wysokości fali, wzdłuż  osi y dla x = 20 (a), mapa (b). Podwojona rozdzielczość czasowa i przestrzenna}
    \label{fig:rozkład_2}
  \end{figure}

  Wykres 3.(a) jest gładszy niż 2.(a), na szczególną uwage zasługuje porównanie rozkład maksymalnego wychylenia do 50 metra. 
  Widać że na wykresie 3.(a) maksymalne wychylenie rośnie, a na wykrsie 2.(a) maleje. 
  Wynika to najprawdopodobniej z błędów numerycznych które inaczej oddziałują na równania w lepszej rozdzielczości.
  Same wychylenia na wzniesieiach również uległy zmianie, na pierwszym wzniesieniu wynosi 0.064408 metrów, na drugim 0.058783 m. Różnica względla to 8.73\%.


\section{Wioski}
Maksymalne wychylenie fali zmiejsza się przy zadanej batymerii. Różnica względa wysokości może się zmiejszyć nawet o 19.91\%. 
Podwojenie rozdzielczości symulacji wykazało zmianę tej różnicy do 8.73\%. Dzięki krztałcie dna pokazanym na rysunku 1, możliwe jest zmiejszenie zagrożenia przed falami. 
W dalszych badaniach można sprawdzic wpływ wysokości wzniesień na rozkład wysokości fali.  



\end{document}